\section{The objects being compared}
\label{sec:objects}

The comparison in this paper is between two balance sheets that arrived at the end
of 2022 holding similar problems on the asset side and entirely different ones on
the liability side. SVB Financial Group and the Charles Schwab Corporation both
carried large unrealised losses on securities classified as held to maturity, both
had accumulated those securities during a period of historically low rates, and
both were marked by the same move in the curve over 2022. One failed in March 2023.
The other did not. The pair is therefore useful for a specific reason: it separates
the size of the asset-side loss from the structure of the funding that has to
sustain it.

\subsection{What each bank disclosed}

The two banks are not symmetric in what they told the public about their own
interest-rate exposure, and the asymmetry is the empirical premise of this paper
rather than an incidental feature of the data.

SVB's annual report for 2021 contains a full economic-value-of-equity profile:
a base level and the effect of four parallel shocks, in both dollars and
percentages \citep{SVB10K2021}. Table~\ref{tab:disclosed} reproduces it. The
report for 2022, filed on 24 February 2023 and two weeks before the bank was
closed, contains net-interest-income sensitivity at the same four shocks and no
economic-value figure at any of them \citep{SVB10K2022}. Schwab disclosed
net-interest-income sensitivity at $\pm$100 basis points in both years and no
economic-value figure in either \citep{Schwab10K2021, Schwab10K2022}. Its
disclosed upside sensitivity fell sharply across the two years, from 14.1\% at
$+100$ basis points in 2021 to 3.6\% in 2022, which is consistent with a book
that had already repriced.

\begin{table}[htbp]
\centering
\caption{Disclosed interest-rate sensitivity. SVB's 2021 report is the only
observation in this sample that carries an economic-value-of-equity profile.
Dashes mark figures the filing does not contain; they are absences in the source,
not omissions here.}
\label{tab:disclosed}
\small
\begin{tabular}{lrrrrr}
\toprule
 & \multicolumn{3}{c}{SVB, 31 Dec 2021} & SVB & Schwab \\
\cmidrule(lr){2-4} \cmidrule(lr){5-5} \cmidrule(lr){6-6}
Shock & EVE (\$m) & $\Delta$EVE (\$m) & $\Delta$EVE (\%) & $\Delta$NII (\%) & $\Delta$NII (\%) \\
 & & & & 31 Dec 2022 & 31 Dec 2022 \\
\midrule
$+200$bp & 14,950 & $-5{,}722$ & $-27.7$ & $3.5$ & -- \\
$+100$bp & 17,799 & $-2{,}873$ & $-13.9$ & $1.8$ & $3.6$ \\
Base     & 20,672 & 0          & 0       & --    & -- \\
$-100$bp & 21,904 & $1{,}232$  & $6.0$   & $-1.8$ & $-3.2$ \\
$-200$bp & 21,308 & $636$      & $3.1$   & $-5.8$ & -- \\
\bottomrule
\end{tabular}
\end{table}

Two features of the disclosed profile matter later. The first is that it is
non-monotonic in the downward direction: a 100 basis point fall adds 6.0\% to
economic value, but a 200 basis point fall adds only 3.1\%. The second is that
the exposure SVB reported was already large. A 27.7\% fall in economic value of
equity at a 200 basis point parallel rise is well beyond the 15\% of Tier~1
capital that the Basel Committee's standard identifies as the threshold for
supervisory attention \citep{BCBS2016}. This was disclosed, in a public filing,
fifteen months before the bank failed.

\subsection{The balance sheets behind those disclosures}

Table~\ref{tab:objects} sets out the quantities that drive the rest of the paper.
Every figure is taken from the filings named above.

\begin{table}[htbp]
\centering
\caption{Funding structure and securities losses at both year-ends, in \$ millions
except where noted. The after-tax taint is the loss that would be recognised if
held-to-maturity securities were sold and the remainder of the portfolio
reclassified; it is shown against book equity to indicate scale, not as a
prediction that it would be incurred.}
\label{tab:objects}
\small
\begin{tabular}{lrrrr}
\toprule
 & \multicolumn{2}{c}{SVB} & \multicolumn{2}{c}{Schwab} \\
\cmidrule(lr){2-3} \cmidrule(lr){4-5}
 & 2021 & 2022 & 2021 & 2022 \\
\midrule
Uninsured deposits          & 182,100 & 165,400 & 134,800 & 68,900 \\
Uninsured share (\%)        & 96.2    & 95.5    & 30.4    & 18.8 \\
Book equity                 & 16,236  & 16,004  & 56,261  & 36,608 \\
Unrealised loss, HTM        & 968     & 15,152  & 0       & 14,138 \\
After-tax taint             & 765     & 11,970  & 0       & 11,169 \\
Taint as \% of book equity  & 4.7     & 74.8    & 0.0     & 30.5 \\
\bottomrule
\end{tabular}
\end{table}

The asset-side problem is comparable in absolute terms. SVB's unrealised
held-to-maturity loss at the end of 2022 was \$15,152m against Schwab's
\$14,138m. Measured against book equity the two diverge, because Schwab's equity
base is larger: the after-tax consequence of recognising those losses would have
consumed 74.8\% of SVB's book equity and 30.5\% of Schwab's.

The funding side is where the two banks are barely the same kind of institution.
At the end of 2022, 95.5\% of SVB's deposits were uninsured, against 18.8\% of
Schwab's. The mechanism connecting the two columns is the one set out by
\citet{Jiang2024}: an unrealised loss becomes a solvency problem only to the
extent that the funding base has both the incentive and the ability to leave
before the loss is absorbed, and uninsured depositors are the ones with the
incentive. The pairing used here is a sharp case of that argument, since the two
banks sit at opposite ends of the uninsured-leverage distribution while holding
similar losses.

\subsection{An accounting difference that has to be removed first}

The two securities portfolios cannot be compared at face value, because Schwab's
disclosed cost basis was reset during 2022. In two transfers, in January and
November, Schwab moved securities with an original cost of \$206,828m out of
available-for-sale and into held-to-maturity at a fair value of \$188,600m,
leaving \$18,228m of loss in accumulated other comprehensive income to be
amortised rather than carried against the securities themselves
\citep{Schwab10K2022}. By the end of the year \$707m had been amortised back,
leaving \$17,521m outstanding.

The effect on any price comparison is large. Measured against the cost basis the
filing reports, Schwab's securities book was worth 92.07 per \$100. Measured
against what the bank originally paid, the same book was worth 87.47 per \$100.
The difference is an accounting election, not an economic quantity, and a
comparison that ignores it measures the election.

SVB is not entirely free of this adjustment. It redesignated \$8,800m of
securities during 2021, carrying \$132m of unrealised loss at the transfer date
\citep{SVB10K2021}. The amount is small enough to leave its cost basis
effectively intact for the purposes of this paper -- \$132m against Schwab's
\$18,228m -- and by the 2022 valuation date no unamortised transfer loss remained
on SVB's books. The comparison in Section~\ref{sec:comparison} is therefore made
against Schwab's original cost rather than its disclosed basis, and the
difference between the two is reported rather than absorbed.
